\section{层次化模仿与技能发现}

\subsection{从演示中提取技能}

模仿学习不再依赖 reward，而是从 demo 中推断 skill segmentation。常见做法包括 sliding window clustering、change point detection、以及 skill-specific autoencoder。

\begin{importantbox}{技能发现 pipeline}
\begin{enumerate}
    \item 打标签：把 demo 分割成 candidate skill；
    \item 训练 skill policy：把 segment 作为 policy context；
    \item 训练 scheduler：高层学习预测 skill sequence；
    \item Deployment：把 skill library 绑定 LLN/LLM 规划器。
\end{enumerate}
\end{importantbox}

\subsection{模仿数据治理}

\begin{knowledgebox}{治理模仿数据的 checklist}
\begin{itemize}
    \item Demo provenance：记录来源视频/轨迹 id；
    \item Skill annotation：标注子任务边界并导出 `labels.json`；
    \item Diversity filter：剔除过度重复的演示片段；
    \item Success flag：每段演示标记是否完成，供 low-layer reward 使用。
\end{itemize}
\end{knowledgebox}

\subsection{VQ-VAE 与离散潜空间}

VQ-VAE 负责把 variable-length 轨迹压缩为 discrete codebook entry，每个 entry 对应一个技能。高层策略只需预测 codebook index，低层 decoder 即可“解码”出需要的动作序列。

\begin{figure}[H]
\centering
\includegraphics[width=0.85\textwidth]{slides-images/slide-25.jpg}
\caption{Skill Discovery 的架构示意（slide-25）。}
\end{figure}

\begin{warningbox}{潜变量 collapse}
若 codebook size 太小，VQ-VAE 会反复复用同一技能，导致高层策略缺乏多样性。常见解决方式是增加 commitment loss 或使用 EM-style clustering。
\end{warningbox}

\subsection{语言作为技能与命令}

LLM 可以充当高层 policy 的语言 interface，讲者提到 SayCan/Inner Monologue 通过 language → intent → skills pipeline 让低层执行器可解释性更高。

\begin{knowledgebox}{语言指令治理要点}
\begin{itemize}
    \item Prompt log 记录 language input + expected skill；
    \item 使用 parser 把 language 转换成 discrete skill id；
    \item 低层反馈（例如 success flag）需要回写到 prompt log，形成 governance loop。
\end{itemize}
\end{knowledgebox}

\subsection{本章小结}

模仿模块让我们可以从演示中自动发现技能库，配合语言/LMM 接口则实现人机可解释的层次调度。

